\chapter{\rm\bfseries Discussion}
\label{ch:discussion}

The results of Chapter~\ref{ch:results} admit a compact reading. A
single mechanical picture --- the progressive engagement of a compliant
blade with the surface --- accounts for the shared force-depth curve,
for the depth-gating of tilt robustness, and for the seat-mount
equivalence (Section~\ref{sec:bell-plus-valley}). The various
\emph{asymmetries} in the record, far from being noise, act as
diagnostics: their sign structure separates what is intrinsic to the
finger from what is contributed by the cloth, the fixture, and the
force estimator (Section~\ref{sec:asymmetries}). The structure of the
measurement noise itself carries methodological lessons for anyone
benchmarking compliant grippers (Section~\ref{sec:manufacturing-mechanism}),
and the three threads together yield concrete design guidance for
cloth-handling systems (Section~\ref{sec:deployment-envelope}).

\section{One Mechanism: How a Compliant Blade Engages Cloth}
\label{sec:bell-plus-valley}
\label{sec:depth-gating-mechanism}
\label{sec:mount-tilt-mechanism}

The bell-plus-valley force-depth shape observed on all four finray
variants (Section~\ref{sec:shape-results}) is consistent with the
finray finger operating in three sequential contact regimes as press
depth increases. In the boundary regime, only the distal tip of the
finger contacts the cloth; force rises roughly linearly with penetration
because the contact area grows in proportion to displacement. At the
knee, the full ridged section of the finger is engaged and the finger
resists compression through its cross-web structure, producing the peak
force. Past the peak, the ribs partially buckle and the finger folds
back on itself; contact area increases but effective stiffness drops,
producing the valley. At recovery depths, the finger has fully wrapped
and further penetration is resisted by structural compression through
the finger base, producing the second rise. That the four features
(boundary, knee peak, valley, recovery) occur at the same calibrated
depths ($\pm 1$\,mm) across all four variants indicates the three-regime
response is a geometric property of the finray shape itself, not
sensitive to finger thickness ($18$ vs $26$\,mm) or to a $5^{\circ}$
mount tilt: force magnitudes scale with contact area, but the depth
locations of the transitions do not.

The same engagement story explains why tilt robustness is depth-gated.
The bipolar bias trajectory of Section~\ref{sec:bipolar-bias-results}
--- a strong signed asymmetry at boundary depths that attenuates to
under $0.9$\,N beyond $d = +4$\,mm --- tracks the contact-geometry
transition documented in Figure~\ref{fig:press-deform}. At boundary
depths only the blade tips touch the surface: the contact patch is
small, and a $2.5^{\circ}$ orientation change redistributes which tip
carries the load, so the asymmetric ridge profile of the finger
dominates the force response. Past the buckling knee, the blades wrap
the surface and the contact becomes an extended line along each blade;
the same orientation change then only slightly re-distributes load
along an already-engaged contact, and the force response becomes
insensitive to tilt sign. The depth at which the bias attenuates
($+4$\,mm) sits just past the knee-peak region
(Section~\ref{sec:shape-results}), consistent with wrap-around
engagement, rather than press depth per se, being the quantity that
confers tilt robustness. This is the mechanism behind the depth-gated
tilt robustness identified as the principal empirical contribution of
this thesis (Chapter~\ref{ch:conclusion}): compliance does not remove
tilt sensitivity, it relocates it to the contact-boundary regime that a
sufficiently deep press avoids.

The mount-tilt equivalence closes the loop on this picture by locating
\emph{where} on the blade the engagement happens. Raising the seat
angle from the $\alpha = +2.5^{\circ}$ baseline to
$\alpha = +5.0^{\circ}$ proved equivalent to a $-1.13$\,mm depth shift,
axis-independent to within $0.02$\,mm
(Table~\ref{tab:mount-tilt-alignment}). Because the comparison variant
(\texttt{finray\_26}) already sits on the $+2.5^{\circ}$ baseline seat
(Section~\ref{sec:variants}), the differential tilt between the two
mounts is $2.5^{\circ}$; a differential tilt of a rigid finger of
length $L$ produces a projected vertical offset at the contact of
$L\left[\sin(5^{\circ}) - \sin(2.5^{\circ})\right] = 0.0436 L$, and
setting this equal to the observed shift gives
$L_{\mathrm{eff}} \approx 26$\,mm --- near the mid-length of the
$\sim 50$\,mm finger, consistent with load distributed along the
finray's curved contact surface rather than concentrated at the distal
tip. The axis-independence is itself a strong constraint: it rules out
any mechanism in which the pre-tilt produces a different effective
load-bearing length on the two axes, so the seat-angle increment acts
as a strictly rigid translation of the operating point in depth,
independent of the runtime tilt axis.

\section{Reading the Asymmetries}
\label{sec:asymmetries}
\label{sec:cross-axis-mechanism}

The engagement mechanism also reframes what the measured tilt
asymmetries \emph{are}. The bias curves of
Sections~\ref{sec:bipolar-bias-results}--\ref{sec:cross-axis-results}
compare the two tilt signs at equal raw command depth, whereas the
per-series calibrated view of Section~\ref{sec:tilt-similarity-results}
--- each orientation zeroed at its own first successful cell --- shows
the five orientation curves nearly coinciding, with mean spreads of
only $1.3$--$2.4$\,N. Taken together, these two views imply that the
raw-axis asymmetry is, to first order, an \emph{orientation-dependent
shift in engagement onset} seen through the local slope of the bell
curve, rather than a change in the force law itself. This reading
predicts the bipolar bias trajectory directly: a bias proportional to
(onset shift)~$\times$~(local slope) is large where the curve is steep
(the boundary), crosses zero near the knee peak where the slope changes
sign --- matching the measured zero crossing between $d = -0.5$ and
$+1.5$\,mm --- and attenuates toward zero in the flat plateau. The
asymmetry nonetheless remains operationally decisive, because a
deployed system calibrates once, blind to its own tilt error, and
therefore lives on the raw axis: the onset shift is exactly the
placement error such a system experiences at light press.

What makes the asymmetries scientifically useful is that their sign
structure points at three different origins. The y-axis bias is
intrinsic to the finger: the finray has mirror symmetry about the
x-axis (the mirror plane between the two opposing fingers) but not
about the y-axis (the gripping direction, which breaks left-right
equivalence by construction), so a $+2.5^{\circ}$ rotation about the
gripping axis loads the two fingers differently and produces a signed
force difference. The x-axis bias, by contrast, should vanish by
symmetry if the finger and cloth were both symmetric about the mirror
plane; that it is nonzero \emph{and of opposite sign to the y-bias}
indicates the symmetry-breaking is not on the finger side. The most
parsimonious explanation is a directional weave-warp axis in the cloth
specimen, misaligned with the gripper's mirror plane, coupling through
the whole-hand contact into a signed force difference --- and the
reproducibility of the x-bias across two independently mounted variants
(\texttt{finray\_18} and \texttt{finray\_26\_5deg}) is consistent with
a shared external source (cloth or table) rather than a per-print
effect.

The third asymmetry separates the two prints of the 18\,mm finger. As
established in Section~\ref{sec:manufacturing-results}, this is not a
categorical behavioural difference: both prints' bias curves share the
same depth-gated envelope --- largest near the boundary, relaxing
toward zero at depth --- and differ only in a shrinking,
print-dependent offset ($2$--$3.5$\,N near the boundary, $<1$\,N by
$d=+10$\,mm) that happens to carry \texttt{finray\_18} across zero
while \texttt{finray\_18\_model2} remains on one side throughout. The
offset lives entirely on the secondary (off-axis) force component, is
comparable to the print-to-print noise floor, and leaves the dominant
$F_z$ response essentially unchanged (shape NMSE $0.41\%$). Two
candidate mechanisms are consistent with the data: small print-instance
differences in rib-wall thickness or infill density manifesting as a
direction-dependent buckling offset; or a small, direction-dependent
bias in the $F_z$-dominant external-wrench estimate
(Section~\ref{sec:frames}), which has documented pose-dependent bias on
other Franka platforms~\cite{petrea2021interaction} and would naturally
be largest where the signal is weakest (the boundary) and vanish where
it is strongest (the plateau).

Each origin comes with its own discriminating experiment. Rotating the
cloth specimen $90^{\circ}$ in the lab frame should flip the sign of
the x-bias while leaving the y-bias approximately unchanged, isolating
the cloth contribution. Distinguishing manufacturing variation from
estimator bias in the print offset requires post-hoc CT-scanning of
the two prints or the independent force-sensor cross-check identified
as future work (Chapter~\ref{ch:conclusion}); until then, the print
offset is a bound on secondary-axis force uncertainty in this dataset,
not evidence that the two prints respond differently to tilt.

\section{Measurement Lessons}
\label{sec:manufacturing-mechanism}
\label{sec:failure-mode-discussion}

Benchmarking compliant grippers on cloth turned out to impose its own
methodology, and two lessons are worth stating explicitly.

The first concerns the noise hierarchy. Same-day print-to-print
variability ($\approx 4.6$\,N mean, RMSE $5.4$\,N) proved
\emph{smaller} than same-print cross-session drift ($7$--$10$\,N).
From a design-decision perspective, any cross-variant force-magnitude
claim with an effect size below the $\approx 5$\,N print-to-print
offset is unlikely to reproduce across different physical instances of
the same CAD, and should be regarded as tentative until replicated
across prints. From an experimental-methodology perspective, session
drift is the dominant noise source in cloth-grasping benchmarks of
this kind: any campaign that pools data across days must either report
cross-day drift as a separate variance component or restrict its
headline claims to same-session pairs. In this thesis, all bias values
are reported for same-session pairs unless explicitly labelled
otherwise; cross-day comparisons are used only for shape-consistency
observations, not for headline effect sizes.

The second lesson concerns what counts as a failure. The trial-outcome
taxonomy of success, no-grasp, and failed trial
(Section~\ref{sec:trial-protocol}) is deliberately three-valued because
the categories map to physically distinct regimes with distinct control
implications. Treating no-grasp trials as failed trials would
systematically delete data from the boundary regime, where cloth
engagement is intermittent but force measurement is clean ---
truncating the observed force-depth curves and masking the
boundary-tilt biases that are one of the headline findings. The
opposite error, treating failed trials as no-grasps, would pollute the
force averages with spurious near-zero readings from aborted
controllers, inflating apparent variance. The three-way taxonomy is the
minimum classification that separates ``the trial produced valid
measurement'' from ``the grasp was successful'' from ``the trial was
invalid.'' It also has a deployment payoff: the no-grasp regime is
precisely where a force-feedback grasp controller could intervene
productively --- contact force is measured and depth is known, but the
lift is failing --- so a closed-loop controller monitoring $|F_z|$
during lift could detect the impending slip and either deepen the
press or abort before dropping the cloth outside the pick zone.
Distinguishing no-grasp from failed-trial in the data record makes
this control strategy visible in the offline analysis.

\section{Implications for System Design}
\label{sec:deployment-envelope}

The parallel-jaw versus finray comparison
(Section~\ref{sec:pj-vs-finray-results}) puts at least a
$4\times$--$7\times$ factor between the two classes in tolerance to
depth command noise --- a lower bound, the finray windows being
right-censored at the analysis-window end. This factor is large enough
that any cloth-handling system whose depth uncertainty exceeds $1$\,mm
should default to a finray-class gripper unless there is a specific
force-precision requirement that the finray cannot meet.

The absence of excess-force safety events on any finray variant across
the full $-2$ to $+16$\,mm range indicates that finray-based systems
can be deployed open-loop with no depth-clipping controller and no
safety-limit tuning. Parallel-jaw systems require an active depth
governor whose accuracy budget must be at least a factor of two
tighter than the successful-grasp window, i.e.\ sub-half-millimetre
command precision, to avoid frequent aborts.

Two further design levers follow from the mechanism sections. The
seat angle offers depth-envelope tuning at essentially no cost: the
finray's operating envelope can be shifted a chosen $1.13$\,mm in
depth per $2.5^{\circ}$ of seat increment without altering the
underlying force-depth sensitivity, at the cost of a $\sim 3$\,N
vertical force offset. And the expanded dataset supports a
three-position design envelope --- safe-plateau finray, boundary finray
with tilt control, and rigid jaw with compliant pad --- with one
refinement: the compliant-pad variant (\texttt{parallel\_jaw\_TPU})
has narrower successful-grasp coverage than its compliant facing might
suggest, and its excess-force susceptibility is closer to
\texttt{parallel\_jaw\_stock} than to finray. The compliant pad
remains a useful hybrid, but its robustness benefit is qualitative
(better boundary behaviour) rather than transforming the parallel jaw
into a finray substitute.

\section{Limitations}
\label{sec:limitations}

\textbf{Single cloth specimen and single surface.} The reported working
ranges and cloth-fixture asymmetries were measured on one cotton cloth
specimen on one acrylic table surface. The absolute numbers will shift
with cloth thickness, weave direction, friction, and surface compliance.
The specific cross-axis sign reversal attributed to cloth weave in
Section~\ref{sec:asymmetries} requires a $90^{\circ}$-rotated
cloth control to isolate the cloth contribution from the fixture
contribution.

\textbf{Cross-session drift.} The dominant noise source in this
dataset is cross-session drift on the same specimen, exceeding
print-to-print manufacturing variability. All same-day comparisons
reported here are internally valid; cross-day claims are qualified
throughout.

\textbf{Parallel-jaw x-axis coverage.} The parallel-jaw variants were
not swept across the x-axis tilt sign; the parallel-jaw bias values
are restricted to the y-axis. Whether the x-axis bipolar bias observed
on the finrays has an analogue on the rigid grippers is untested.

\textbf{Boundary-depth sample sizes.} Boundary-depth cells frequently
have $n=3$ and inter-trial standard deviations of $0.5$--$1.5$\,N
driven by cloth-fold state variability that is not further
characterised. Boundary values are honest measurements with
correspondingly wide error bars.

\textbf{Success rate as a binary rating.} The success/no-grasp
distinction is operator-assessed after the lift stage; a
finer-grained grip-force-during-lift measurement would resolve
marginal cases that the current binary rating collapses.

\section{Code and Data Availability}

All source code, raw trial data, ground-truth JSON files, and analysis
scripts are published in the project repository at
\TODO{public GitHub URL}. The thesis was written against commit
\TODO{commit hash}.
